\section{Local shielding of singular sets}\label[appendix]{sec:local}

Shielding makes each singular point irrelevant to the next curvature
requirement by enlarging enough remaining pole cores near that point.
The forcing costs a potential concentrated in a thin annulus.  We
first pay for this potential from the reserved gradient form; the
next appendix makes the choices compatible throughout the diagram.

We work at one fixed stage with a finite set of remaining labels.
The countable extension is proved in \cref{countable:step}.  The local
estimates use only $.60\le a_r\le .63$ and the positive reserve of
\cref{cut:reserve}.  Smooth geometric bounds are taken on the specified
compact windows, so no bounded-geometry assumption on the original
manifold is involved.

Fix a true cut in a parent packet $E$ of base dimension $m$, as in
\cref{cut:gmt}.  Its boundary support is $S$, its singular set is $\Sigma$,
and its regular part is $S_{\rm reg}$.  The old torus volume factor is
$\omega$.  All measures below include this factor; on a cylinder it is
pulled back from $S_{\rm reg}$.  For either phase $\Omega$, write
\[
 J_\Omega=\overline\Omega_{\rm reg}\mathbin{\cup}_{S_{\rm reg}}
             (S_{\rm reg}\times[0,\infty)).
\]
Here the notation retains only regular boundary points when forming the
join.  Tests on $J_\Omega$ are compactly supported, piecewise smooth
functions with matching values at the seam, whose supports avoid
$\Sigma$ and the cylinders above it.  Lipschitz tests with the same
properties follow by approximation in regular charts.  Tests on the
wall have the analogous meaning.  Use the reserved energy from
\cref{cut:reserve}, namely
\begin{equation}\label{local:norm}
 \mathcal N(v)=
 \begin{cases}
  1.7\displaystyle\int_{S_{\rm reg}}|\nabla v|^2,&\text{on the wall},\\[3pt]
  .8\displaystyle\int_\Omega|\nabla v|^2+
  1.7\displaystyle\int_{S_{\rm reg}\times[0,\infty)}|\nabla v|^2,
       &\text{on a joined side}.
 \end{cases}
\end{equation}
The notation $\|v\|_2^2$ denotes the sum of the corresponding weighted
$L^2$ integrals.  An ambient function is restricted to the wall and
frozen in the cylinder coordinate.  In particular, these conventions
specify every potential integral in this section.  The true cut remains
fixed throughout: only profiles of labels other than the processed
label $j_0$ will change.

\begin{lemma}[Choice of anchors]\label{local:anchors}
There are compatible smooth functions $h_j$, for the remaining labels,
within $.1$ of $\max(-1,\min(1,\mu_j))$, and common numbers
$t_j\in(-.3,.3)$ such that, on the singular set in each parent of
dimension $m$, at most $m-8$ of the equalities $h_j=t_j$ hold.
The assertion is vacuous in dimensions below eight.  Moreover
\[
 \{h_j>t_j\}\cap\{\mu_j=-\infty\}=\varnothing,
 \qquad
 \{h_j<t_j\}\cap\{\mu_j=+\infty\}=\varnothing.
\]
\end{lemma}
\begin{proof}
Compatible smoothing by \cref{packet:approximation} gives the stated
approximations: smooth the continuous clipped functions on compact
parts and blend with their already prepared exact face copies.
Allocate the error $.1$ beforehand over the finitely many requests
on each face, and reproduce constants exactly on passive copies.
For a packet of dimension $m\ge8$ and a set $I$ of $m-7$
remaining labels, the smooth map $h_I=(h_j)_{j\in I}$ is locally
Lipschitz.  Since $\dim_{\rm H}\Sigma\le m-8$, its image of $\Sigma$
has zero $(m-7)$-dimensional Lebesgue measure.  This follows on each
relatively compact chart from the defining Hausdorff covers and then
on the packet by a countable union.  The set of threshold vectors for
which $t_I\in h_I(\Sigma)$ is null by Fubini.  For measurability,
intersect $\Sigma$ with a compact exhaustion first: each image is compact,
and their countable union is measurable.  Give each interval $(-.3,.3)$
its normalized Lebesgue measure.  There are countably many packets and
finitely many relevant subsets per packet, so choose a threshold vector
outside the union of all exclusions.  The same argument also works in
a countable product of the threshold intervals: each exclusion is a
null cylinder set depending on only $m-7$ coordinates.  If fewer
than $m-7$ labels remain, the asserted count holds without a condition.
At a negative pole $h_j\le-.9<t_j$, and at a positive pole
$h_j\ge.9>t_j$, proving the last assertions.
\end{proof}

The guard $H_E^*$ in \eqref{packet:guard} requires at least $m-2$
remaining undecided labels.  At $x\in\Sigma$, enlarging the positive
core for every label with $h_j(x)>t_j$ and the negative core for every
label with $h_j(x)<t_j$ to contain $x$ would leave at most
$m-8<m-2$ labels undecided.  These pointwise decisions would therefore
remove $x$ from the guard.

We need quantitative pieces of $\Sigma$, rather than just its
Hausdorff dimension.  A coordinate chart in the following lemma is
chosen with compact buffer and with its metric as close to Euclidean
as required.  Such charts can be chosen around every point.  The plane
condition is expressed in Euclidean chart coordinates.  Replacing
Euclidean balls by metric balls changes only fixed factors, which can
be absorbed by requesting a smaller flatness constant.

\begin{lemma}[Compact pieces with one-sided flatness]\label{local:flat-pieces}
Let $O$ be an open subset of a parent packet and let $\eta>0$.
The set $\Sigma\cap O$ is covered by countably many compact sets
$F\subset\Sigma\cap O$, each contained in a buffered coordinate chart,
with the following property.  There is $r_F>0$ such that, for every
$x\in F$ and $0<r<r_F$, some affine $(m-8)$-plane $L_{x,r}$ satisfies
\begin{equation}\label{local:flatness}
 F\cap B_r(x)\subset\{z:\dist_{\rm Eucl}(z,L_{x,r})<\eta r\}.
\end{equation}
The plane may depend on both the center and the radius.
\end{lemma}
\begin{proof}
Write $\Theta(x)$ for the unweighted boundary density.  First consider
a tangent minimizing boundary cone $C$ at a singular point.  It is
nonflat, since a flat tangent gives regularity.  For any $z\in C$, the
density ratio centered at $z$ has limit $\Theta_C(0)$ as its radius
tends to infinity: balls centered at $z$ are sandwiched between balls
centered at zero whose radii differ by $|z|$.  Consequently
\[
 \Theta_C(z)\ge\Theta_C(0)
 \quad\Longrightarrow\quad
 \text{every density ratio centered at $z$ equals $\Theta_C(0)$}.
\]
Equality in monotonicity says that $C$ is also a cone centered at $z$.
Composing a dilation about $z$ with the inverse dilation about zero
shows translation invariance along $z$.  Conversely, translation
invariance along $z$ gives equality of the two densities.  The
translation invariance directions form a vector subspace, the spine.
Every spine point is singular: otherwise translation gives regularity
at the vertex, and the cone is a plane.  The minimizing-cone singular
dimension estimate therefore bounds the spine dimension by $m-8$.
These tangent-cone and regularity inputs are precisely those in
\cref{cut:gmt}.

Fix a singular point $x$ and a small number $\eta'>0$.  There are
$d_x>0$ and $r_x>0$ such that, at every radius $r<r_x$, the points
\begin{equation}\label{local:high-density}
 \{z\in\Sigma\cap B_r(x):\Theta(z)\ge\Theta(x)-d_x\}
\end{equation}
lie within $\eta'r$ of some $(m-8)$-plane through $x$.  If not, choose
failing radii $r_i\downarrow0$ with density tolerances tending to zero.
Rescale at $x$ and pass to a tangent cone $C$.  Fix a plane containing
its spine.  Failure for this particular plane gives points $z_i$ in
the rescaled unit ball a fixed positive distance from it.  A convergent
subsequence has limit $z$.  Upper semicontinuity of densities under
these convergent blowups gives
$\Theta_C(z)\ge\Theta(x)=\Theta_C(0)$.  The preceding paragraph puts
$z$ in the spine, a contradiction.  Smooth metric errors and bounded
mean-curvature errors vanish in this rescaling.

Here is the countable decomposition, including its closure step.
Group singular points according to lower bounds
$r_x\ge1/a$ and $d_x\ge1/b$, with positive integers $a,b$, and according
to a density interval of length less than $1/(2b)$.  Any two points
$x,z$ in the same group satisfy $\Theta(z)>\Theta(x)-d_x$.
Thus the group itself has the plane containment at every one of its
centers and every radius less than $1/a$.  Work within one buffered chart at a time.  Intersect these groups
with a countable family of compact subpatches $Q\Subset O$ exhausting
that chart's portion of $O$, and take
$F=\overline{Q\cap\text{group}}$.  Relative closedness of $\Sigma$
and the compact buffer ensure $F\subset\Sigma\cap O$.
These compact sets still cover $\Sigma\cap O$.

To check the geometric condition at a closure center $x$, let
$x_i$ be group points tending to $x$ and fix $r<1/(3a)$.  Apply the
group estimate in $B_{2r}(x_i)$.  The planes pass through $x_i$, so a
subsequence converges to a plane through $x$.  Every group point in
$B_r(x)$ is in these larger balls eventually.  Passing first to the
limit of the planes and then to closure points gives containment with
width $2\eta'r$.  Taking, for example, $\eta'<\eta/3$ proves
\eqref{local:flatness}; choose $r_F$ below $1/(3a)$ and the buffered
chart radius.  This argument transfers the containment
itself; no density continuity at a closure point is required.
\end{proof}

The next construction fixes the pressure and potential to be estimated.
Distances denoted by $\delta$ from now on are ambient base distances.

\begin{lemma}[A shell forcing profile]\label{local:shell}
Let $F$ be a compact subset of $\Sigma\cap\{h_j>t_j\}$.
For every sufficiently small $\varepsilon>0$ there is an extended-real
profile $u_\varepsilon$, equal to $+\infty$ for $\delta\le\varepsilon$
and $-\infty$ for $\delta\ge R_\varepsilon$, where
\[
 \delta(x)=\dist(x,F),\qquad R_\varepsilon=e^{5\pi}\varepsilon,
\]
whose finite truncations are locally Lipschitz and which satisfies
\begin{equation}\label{local:shell-gradient}
 2|\nabla u_\varepsilon|
       \le .58u_\varepsilon^2+2\delta^{-2}
       \quad (\varepsilon<\delta<R_\varepsilon).
\end{equation}
There is a smooth nonnegative ambient function $D_\varepsilon$ with
\begin{equation}\label{local:shell-potential}
 \begin{aligned}
 D_\varepsilon&\ge2\delta^{-2}
       &&\text{on }\{\varepsilon\le\delta\le R_\varepsilon\},\\
 D_\varepsilon&\le2.12\delta^{-2}
       &&\text{everywhere off }F,\\
 \supp D_\varepsilon&\subset
       \{\varepsilon/2<\delta<2R_\varepsilon\}
       \Subset\{h_j>t_j\}.
 \end{aligned}
\end{equation}
Replacing $\mu_j$ by $\max(\mu_j,u_\varepsilon)$, adding
$D_\varepsilon$ to $\Pi_j$, and adding
$\{\delta<\varepsilon/2\}$ to $U_j^+$ preserves the profile inequality
with coefficient $a_r$.  There is an identical negative construction
in $\{h_j<t_j\}$, using $\min(\mu_j,-u_\varepsilon)$ and enlarging
$U_j^-$.  All supports and newly forced regions are compactly
contained in the stated strict anchor region.
\end{lemma}
\begin{proof}
The polynomial
$q(y)=.29y^2-y+1$ is strictly positive, and
\[
 \int_{-\infty}^{\infty}\frac{dy}{q(y)}=5\pi.
\]
Let $y:(0,5\pi)\longrightarrow\R$ solve $y'=q(y)$, with the two
infinite endpoint limits, and put
\[
 u_\varepsilon(x)=\frac{y(\log(R_\varepsilon/\delta(x)))}{\delta(x)}.
\]
At differentiability points of $\delta$, its gradient has norm at
most one.  Differentiating in $\delta$ gives
\[
 \left|\frac{d}{d\delta}\frac{y(\log(R_\varepsilon/\delta))}{\delta}\right|
   =\frac{y'+y}{\delta^2}
   =.29u_\varepsilon^2+\delta^{-2},
\]
which proves \eqref{local:shell-gradient}.  Extending by the stated
poles is continuous in the extended real line.  Each finite clipping
has a finite Lipschitz constant locally, including across either
clipping interface.

Choose a continuous cutoff of $\delta$, equal to one on
$[\varepsilon,R_\varepsilon]$ and supported strictly in
$(\varepsilon/2,2R_\varepsilon)$.  Multiplying it by
$2.06\delta^{-2}$ gives a continuous nonnegative function supported
away from $F$.  Smooth it by nonnegative local convolutions and a
partition of unity, with sufficiently small value error on this
compact annulus and with support retained in a slightly larger
annulus of the indicated size.  The positive margins $.06$ on the
central annulus and $.06$ below the upper bound give
\eqref{local:shell-potential}.  The annuli lie in the strict anchor
region for all small $\varepsilon$, by compactness of $F$.

At almost every point where the new maximum is finite, its gradient
is the gradient of an active branch.  The old branch obeys the old
inequality.  The new one obeys it because $.58<a_r$ and
$D_\varepsilon\ge2\delta^{-2}$ wherever that branch is finite.
The same reasoning applies at branch equality, using the elementary
weak derivative rule for maxima of locally Lipschitz functions.
The new core has closure inside the open positive pole region
$\{\delta<\varepsilon\}$.  By \cref{local:anchors} no old negative
pole occurs in the support of the positive forcing.  Positive and
negative additions for the same label have disjoint strict anchor
regions, so repetitions introduce no conflict.  Indeed any opposite
pole introduced later lies in the opposite strict anchor region and
therefore cannot meet this support.  Original poles and previously
introduced opposite cores consequently remain protected throughout.
Replacing the profile by its negative proves the other assertion.
\end{proof}

We now prove a Hardy estimate with a positive, dimension-dependent
coefficient.  Its coefficient must be chosen before the more accurate
flatness used later.

\begin{lemma}[Uniform rough Hardy coefficient]\label{local:rough}
For each $m\ge8$ there are numbers $\eta_0>0$ and $c_0>0$ with the
following property.  Suppose $F$ satisfies \eqref{local:flatness}
with $\eta\le\eta_0$ in a sufficiently accurate buffered chart.
There are $d_F>0$ and $C_F<\infty$ such that, on the wall and on
either joined side,
\begin{equation}\label{local:rough-estimate}
 c_0\int_{\delta<d_F}\delta^{-2}v^2
             \le\mathcal N(v)+C_F\|v\|_2^2.
\end{equation}
The constants $\eta_0,c_0$ depend only on $m$ and a fixed permitted
chart distortion.  The constants $d_F,C_F$ may depend on the compact
piece, its scale threshold, and the local smooth metric, weight,
and pressure.  No principal-curvature bound for the wall is assumed.
\end{lemma}
\begin{proof}
Put $s=m-8$.  At scale $r$, the one-sided plane condition covers
$F\cap B_r(x)$ by at most $C_m\lambda^{-s}$ balls of radius
$\lambda r$, with centers in $F$, where $\lambda=8\eta_0$ can be
made as small as desired.  To see the bound, take a maximal
$\lambda r$-separated family in the set.  Projection to the
approximating $s$-plane changes pairwise distances by at most
$2\eta_0r$; the projections are separated by a fixed multiple of
$\lambda r$ in an $s$-ball of radius $2r$.  Its elementary Euclidean
packing bound gives the claim, with harmless changes in $C_m$ and
the covering radius.  Absorb these fixed changes into the initial
choice of $\lambda$.  Choose $\eta_0$ so small that the resulting
covering number is at most $\lambda^{-(s+1/2)}$.

Iterating the same covering at centers in $F$ gives, for $0<t\le d$
below a fixed fraction of $r_F$, a covering by at most
$C(d/t)^{m-7.5}$ balls of radius $t$ of the part of $F$ controlling a
ball of radius $d$.  Here one first centers a ball of radius a fixed
multiple of $d$ at a nearby point of $F$ and then iterates; the
constants use only the fixed rough choice.  Enlarging the final
balls by a fixed factor proves the tubular-volume estimate
\begin{equation}\label{local:tube-volume}
 |B_d(x)\cap\{\delta<t\}|
       \le C d^{m-7.5}t^{7.5},\qquad 0<t\le d,
\end{equation}
whenever $\dist(x,F)\le d$ and the slightly enlarged ball is in the
chart at allowed scales.  Metric and Euclidean distances are
uniformly comparable here.  Dyadically splitting the regions
$2^{-k-1}d<\delta\le2^{-k}d$, and treating $\delta>d$ separately,
gives
\begin{equation}\label{local:density-integral}
 \int_{B_d(x)}\delta^{-5.5}\,dx\le C d^{m-5.5}.
\end{equation}
Indeed the sum of the bounds is a geometric series with ratio
$2^{-2}$.  In particular this density is locally integrable.

Choose a nonnegative smooth chart cutoff $\chi$, equal to one on a
neighborhood of $F$, and, using Euclidean coordinates, define
\begin{equation}\label{local:rough-potential}
 U(x)=1+\int\chi(y)\delta(y)^{-5.5}|x-y|^{-(m-4)}\,dy.
\end{equation}
For small $\delta(x)>0$, the part $|x-y|<c\delta(x)$ has
$\delta(y)$ comparable to $\delta(x)$; it supplies the lower bound
for $U$.  The same region supplies the upper bounds for $U$ and its
first derivatives.  On the successive larger annuli use
\eqref{local:density-integral}; their contributions have orders
$(2^k\delta)^{-1.5}$ and $(2^k\delta)^{-2.5}$ respectively.
The parts above the allowed local scale are bounded constants
depending on $F$.  Shrinking the final distance threshold absorbs
them, and gives dimension-dependent constants $c,C>0$ such that
\begin{equation}\label{local:rough-U-bounds}
 c\delta^{-1.5}\le U\le C\delta^{-1.5},
 \qquad |\nabla U|\le C\delta^{-2.5}.
\end{equation}

Set $a=m-4$ temporarily.  For any Euclidean $d'$-plane $T$, direct
differentiation of $K(z)=|z|^{-a}$ gives
\[
 -\tr_T D^2K(z)
 =a|z|^{-a-2}\left(d'-(a+2)|\operatorname{proj}_T(z/|z|)|^2\right)
 \ge a(d'-a-2)|z|^{-a-2}.
\]
For $d'=m$ or $m-1$ the last coefficient is positive.  A sufficiently
small, dimension-dependent metric distortion preserves a fixed
positive fraction of it for traces in metric orthonormal frames.
The connection and weight terms in the weighted Laplacian involve
only first derivatives of $U$.  On $S_{\rm reg}$ the restriction
formula adds its mean-curvature vector paired with $\nabla U$.
This vector is locally bounded, since
$H_{\rm base}=p-\partial_\nu\log\omega$ by \cref{cut:gmt}.
There is no dependence here on individual principal curvatures.
Integrating the positive second-derivative kernel over
$|x-y|<c\delta(x)$ gives $c\delta^{-3.5}$.  Equations
\eqref{local:rough-U-bounds}, followed by another reduction of the
distance threshold, therefore imply, in the ambient space and on
the regular wall,
\begin{equation}\label{local:U-superharmonic}
 -\frac{\Delta_\omega U}{U}\ge c\delta^{-2},
 \qquad |\nabla\log U|\le C\delta^{-1}.
\end{equation}
The leading constants again depend only on the fixed rough data.
For example the first-order errors before division have size at
most $C_F\delta^{-2.5}$, one power lower than the positive term.
The differentiated kernels through order two are locally integrable
off $F$: their strongest singularity is $|x-y|^{-(m-2)}$.
Since the density is bounded near any such $x$ and integrable
elsewhere on its compact support, ordinary differentiation under
the integral, with a vanishing small diagonal ball, proves that
$U$ is $C^2$ off $F$.  No diagonal distributional term is present.

For a weighted domain $Q$, completing a square and integrating by
parts gives the vector-field identity in inequality form
\begin{equation}\label{local:vector-hardy}
 \int_Q|\nabla v|^2\ge
 \int_Q(\Div_\omega A-|A|^2)v^2
       -\int_{\partial Q}(A\cdot\nu_Q)v^2.
\end{equation}
Choose a smooth cutoff $\rho$, equal to one near $F$ and supported
in the chart, and set $A_b=-b\rho\nabla\log U$ in the bulk.  On the
wall use $A_h=-b\rho\nabla_S\log U$.  For $0<b\le1$, where $\rho=1$,
\[
 \Div_\omega(-b\nabla\log U)-b^2|\nabla\log U|^2
 =-b\Delta_\omega U/U+b(1-b)|\nabla\log U|^2
 \ge bc\delta^{-2}.
\]
Outside a smaller neighborhood of $F$ all errors from these fields
are bounded by a piece-dependent constant.

On the cylinder take the horizontal field $A_h$ and add the vertical
component $a(x)\partial_t$, where
\[
 a(x)=\frac{.8}{1.7}\,b\rho(x)
               \bigl(1+|\nabla\log U|_g^2\bigr)^{1/2}
                 \big|_{S_{\rm reg}}.
\]
It has zero divergence in the vertical direction and satisfies
$1.7a\ge .8|A_b|$ at the seam.  The outward cylinder normal there is
$-\partial_t$.  Applying \eqref{local:vector-hardy} with coefficients
$.8$ and $1.7$ on the two pieces, their boundary terms sum to
\[
 \int_{S_{\rm reg}}
       \bigl[-.8 A_b\cdot\nu_\Omega+1.7a\bigr]v(0)^2\ge0.
\]
This explicitly permits arbitrary matching seam values.  The
vertical square costs at most $C b^2\delta^{-2}+C_F$ in the
cylinder interior.  Choose $b>0$ sufficiently small, using the
universal constants in \eqref{local:U-superharmonic}, so that this
is absorbed by half the positive term of order $b\delta^{-2}$.
All remaining localization terms cost $C_F\|v\|_2^2$.
This proves \eqref{local:rough-estimate} on a joined side.  The
wall proof is the same calculation without the seam or the vertical
field.  Taking the smaller leading coefficient makes one $c_0$
valid for the wall and both sides.
\end{proof}

\begin{proposition}[The local shielding estimate]\label{local:sharp}
There is a sufficiently small flatness number $\eta>0$, depending
only on $m$ and the fixed numerical coefficients, with the following
property.  Let $F$ be any piece from \cref{local:flat-pieces} for this
flatness number, in a sufficiently accurate chart.  For the potential
of \cref{local:shell}, there are $\varepsilon_F>0$ and $C_F<\infty$
such that, for every $0<\varepsilon<\varepsilon_F$,
\begin{equation}\label{local:sharp-estimate}
 1.16\int D_\varepsilon v^2
           \le .78\mathcal N(v)+C_F\|v\|_2^2.
\end{equation}
This holds on the true wall and either joined side, and hence on
translated compact tests in a pure wall cylinder.  The constants
$\varepsilon_F,C_F$ may depend on the local smooth data and $F$,
but are independent of $\varepsilon$.  No vanishing trace condition
is imposed at a joined seam.
\end{proposition}
\begin{proof}
We give the parameter order first.  Fix the rough number $\eta_0$
and coefficient $c_0$ from \cref{local:rough}.  Below, constants in
partition gradient bounds depend only on dimension and the fixed
chart distortion.  Choose $\xi>0$ so small that the rough estimate
pays their total squared-gradient error with at most $.04$ of
$\mathcal N$.  Next choose a large fixed $K$, then a small fixed
$\alpha>0$, and finally a strong flatness $\eta\le\eta_0$, as
specified below.  Only after the piece $F$ has been selected for
this $\eta$ do we choose its small distance threshold and then
$\varepsilon_F$.  In particular $C_F$ is never used to choose
$\xi$, $K$, or $\eta$.

Put $a_-=1/2$, $a_+=2e^{5\pi}$ and $R_1=K\varepsilon$.
The shell supporting $D_\varepsilon$ lies in
$a_-\varepsilon<\delta<a_+\varepsilon$.
Choose a smooth angle $\vartheta(\delta)$ equal to zero for
$\delta\le a_+\varepsilon$ and equal to $\pi/2$ for
$\delta\ge\xi R_1$, with
\[
 |\nabla\vartheta|\le\xi/\delta.
\]
Such a function is obtained by a smooth transition in
$\log(\delta/(a_+\varepsilon))$ of length at least a fixed multiple
of $1/\xi$.  Thus it suffices to choose
$K\ge (a_+/\xi)\exp(C/\xi)$, independently of the piece.  Set
$q=\cos\vartheta$ and $q_\infty=\sin\vartheta$.
They extend as constants across $F$, and
$|\nabla q|^2+|\nabla q_\infty|^2\le\xi^2\delta^{-2}$.

Take a smooth normalized squared partition $\sum_i\phi_i^2=1$
on chart boxes of scale $R_1$, with uniformly bounded overlap and
$\sum_i|\nabla\phi_i|^2\le C_mR_1^{-2}$.
One construction translates a fixed compactly supported bump on a
Euclidean lattice and divides each bump by the square root of the
sum of their squares.  Restrict the same ambient functions to the
wall and freeze them on the cylinder.  Only finitely many boxes
meet $\supp q$, which is compact in the chart at the scales in use.
The family $q\phi_i,q_\infty$ is a squared partition, and its total
gradient square is at most $C_m\xi^2\delta^{-2}$ on
$\{\delta<\xi R_1\}$ and zero elsewhere.  The estimate follows
from $\delta<\xi R_1$ wherever $q\ne0$.

For a squared partition $\sum_i\chi_i^2=1$, direct expansion gives
\begin{equation}\label{local:IMS}
 \sum_i\mathcal N(\chi_iv)
 =\mathcal N(v)+
   \int a(x)\sum_i|\nabla\chi_i|^2v^2,
\end{equation}
where $a(x)$ is $.8$ on bulk and $1.7$ on wall or cylinder.
The mixed terms vanish because $\sum_i\chi_i\nabla\chi_i=0$.
Choose $\varepsilon_F$ small enough that $\xi R_1<d_F$.
Equation \eqref{local:rough-estimate} and the choice of $\xi$ now
give
\begin{equation}\label{local:localization}
 \sum_i\mathcal N(q\phi_iv)
 \le1.04\mathcal N(v)+C_F\|v\|_2^2.
\end{equation}
For example it is enough to impose
$1.7C_m\xi^2/c_0\le.04$; the other partition term is nonnegative
and has been discarded.  Lipschitz distance cutoffs are allowed by
approximation on the regular test supports.

Consider one box meeting the support of $q$.  Choose a nearby point
of $F$ and apply \eqref{local:flatness} at a radius $L_mR_1$,
where a fixed $L_m$ includes the box and every nearest point of
$F$ needed in it.  Taking $\varepsilon_F$ smaller ensures this is
below $r_F$.  Let $L$ be the resulting $(m-8)$-plane and define
\[
 \zeta^2=\dist_{\rm Eucl}(x,L)^2+\alpha^2\varepsilon^2,
 \qquad P=\tfrac12\zeta^2.
\]
The Euclidean Hessian of $P$ is an orthogonal projection of rank
eight.  Its trace on every $(m-1)$-dimensional subspace is at least
seven.  Its gradient has size at most $\zeta$, and throughout the
box $\zeta\le C_mR_1$ (after fixing $\alpha$ and $K$).

For either the ambient or the tangential weighted derivatives,
the vector field $A=d_1\nabla\log\zeta$ has potential
\begin{equation}\label{local:plane-potential}
 (\Div_\omega A-|A|^2)\zeta^2
 =d_1\Delta_\omega P
       -(2d_1+d_1^2)|\nabla P|^2/\zeta^2.
\end{equation}
Let $\tau$ bound the chart's metric distortion.  The rank traces
have error at most $C_m\tau$, and
$|\nabla P|^2/\zeta^2\le1+C_m\tau$.  The connection, weight, and
wall mean-curvature contributions to $\Delta_\omega P$ have size
at most $C_FR_1$.  Thus one may arrange
$C_m\tau\le .005$ and $C_FR_1\le .025$, uniformly in the box and
its plane.  Equation \eqref{local:plane-potential} then gives
\[
 \begin{array}{ll}
 d_1=1\text{ in bulk}:&
       \Div_\omega A-|A|^2\ge4.9\zeta^{-2},\\
 d_1=2\text{ on wall}:&
       \Div_\omega A-|A|^2\ge5.9\zeta^{-2}.
 \end{array}
\]
Indeed the respective Euclidean lower bounds are $8-3=5$ and
$2\cdot7-8=6$.  Only the mean curvature of the wall enters this
calculation, so these bounds remain valid arbitrarily near its
singular set on the regular test supports.  The positive term
$\alpha^2\varepsilon^2$ also makes the fields nonsingular along
$L$; the estimates hold inside the inner shell as well.

In the bulk use $A_b=\nabla\log\zeta$.  On a joined cylinder use
$A_c=2\nabla_S\log\zeta+\beta\zeta^{-1}\partial_t$, where
$\beta=(.8/1.7)\sqrt{1+C_m\tau}$ is chosen so that
$\beta/\zeta\ge(.8/1.7)|A_b|$ at the seam.
The metric distortion determines this fixed upper factor.
The vertical component has zero divergence and is orthogonal to
the horizontal component.  Its only interior cost is
$\beta^2\zeta^{-2}$.  The boundary terms in
\eqref{local:vector-hardy} are nonnegative after summing, by exactly
the cancellation used in \cref{local:rough}.  After multiplication
by the coefficients in \eqref{local:norm}, the limiting interior
coefficients are
\[
 .8\cdot5=4\quad\text{in bulk},\qquad
 1.7\cdot6-.8^2/1.7>9.8\quad\text{in the cylinder},
\]
and $1.7\cdot6=10.2$ on the wall alone.  With the explicit error
bounds just imposed, the coefficients are at least $.8\cdot4.9=3.92$
in bulk, $1.7\cdot5.9-(.8^2/1.7)1.005>9.6$ in the cylinder,
and $1.7\cdot5.9>10$ on the wall.  In particular all exceed $3.8$.
Applying the vector inequality to $q\phi_iv$ gives
\begin{equation}\label{local:box-bound}
 \mathcal N(q\phi_iv)\ge
       3.8\int\zeta^{-2}(q\phi_iv)^2.
\end{equation}
This is a bound throughout the localized support; its nonnegative
right side may therefore be restricted to the shell.

We finally compare $\zeta$ with $\delta$ on that shell.  A nearest
point $z\in F$ lies in the ball used for the plane condition, so
\[
 \zeta(x)\le 1.005\delta(x)+L_m\eta K\varepsilon
                                  +\alpha\varepsilon.
\]
The first factor is ensured by the fixed chart accuracy and its
buffer, after shrinking the scale so that a minimizing short path
to $F$ stays in the chart.  Since $\delta\ge a_-\varepsilon$ there,
choose $\alpha/a_-\le .005$ and then
$L_m\eta K/a_-\le .015$.  These choices give
$\zeta\le1.03\delta$.  Notice that $K$ was already fixed using the
rough coefficient, and that the strong flatness was chosen before
selecting $F$.  Since $3.8/(1.03)^2>3.5$, \eqref{local:box-bound}
controls $3.5\delta^{-2}$ on the shell.  Summing over boxes uses
$q=1$ there and $\sum_i\phi_i^2=1$, and then
\eqref{local:localization} gives
\[
 3.5\int_{\supp D_\varepsilon}\delta^{-2}v^2
                  \le1.04\mathcal N(v)+C_F\|v\|_2^2.
\]
Together with \eqref{local:shell-potential}, this proves the stronger
coefficient bound
\[
 1.16\int D_\varepsilon v^2
 \le\frac{1.16\cdot2.12\cdot1.04}{3.5}\mathcal N(v)
                   +C_F\|v\|_2^2,
 \qquad
 \frac{1.16\cdot2.12\cdot1.04}{3.5}<.731<.78.
\]
This yields \eqref{local:sharp-estimate}.  A compact test on a pure
cylinder can be translated to positive height in a joined cylinder
and extended by zero onto the bulk, so it obeys the same estimate.
\end{proof}

\begin{remark}\label{local:choice-order}
The choices can be summarized without a hidden dependence on the
compact piece: dimension and chart distortion; rough flatness
$\eta_0$ and coefficient $c_0$; localization parameter $\xi$;
logarithmic taper and $K$; $\alpha$ and the stronger flatness $\eta$;
the compact pieces supplied for that $\eta$; finally their individual
distance thresholds and shell scales.  Shrinking a threshold absorbs
large-scale kernel tails and smooth geometric errors.  It never
changes the leading coefficient used to select the stronger
flatness.  All estimates are weighted base estimates, so the number
of previously appended circle factors does not replace $m$ in any
dimension or rank count.
\end{remark}
